\section{Finite comparison maps for the compact realization}
\label{app:eight-finite-constructions}

The finite data of Section~\ref{subsec:eight-finite-constructions}
become useful when their maps can be calculated. This appendix fixes
the geometric supports, coefficient morphisms and matrix equations
needed for that calculation. It also separates the formal target
from the compact source: a target degree can be listed before a
sheaf of that degree or a Hall operation has been constructed.

Throughout, \(X=S\times E\), with \(S\) a smooth projective K3
surface and \(E\) a smooth elliptic curve. The stage parameter
\(R\) bounds the chosen geometric data. It is distinct from root
height, word length, cohomological degree and any formal deformation
parameter. When several bounds are used, their comparison is part
of the construction.

The proposed compact datum includes a field complex
\(\mathcal A_{X,R}^{E_3}\), its local/wrapped specialization,
critical coefficients and orientation, Hall operations, primitives,
a coalgebra and its associative comparison, Pfaffian sections and
recognition maps. If an \(E_3\)- or BV description is required,
its formal target charts, field differential and operations must be
constructed. The quantum master equation, anomaly and framing,
compact-support maps and factorization descent are additional
conditions on those same complexes. A map denoted
\(\operatorname{Sp}_{K3,E}^{\mathrm{ch}}\) must preserve the
specified operations and stage transitions; it is not defined by
naming a current-envelope target.

\subsection{Charges, finite supports and target degrees}
\label{app:build-charge-window}

Fix the Abramovich--Polishchuk heart and its Liu tilt at the scope
of Appendix~\ref{app:stability-support}
\cite{AbramovichPolishchukTStructures,LiuProductStability}.
The polynomial charge has the form
\[
 L_M(n)=a(M)n+b(M)+i\bigl(c(M)n+d(M)\bigr).
\]
The weak charge, torsion pair and tilted heart must be kept with
the central charge; retaining only the polynomial does not retain
the stability condition. A finite subobject- and quotient-closed
collection with a strict integer length admits the finite HN
argument of Proposition~\ref{prop:retained-window-hn-filtration}:
choose maximal phase and then maximal length at each quotient.
Its factor count is bounded by the initial length. This is an
argument in the chosen heart, not a theorem that the moduli of all
such factors is finite or proper.

Let \(\Gamma_{\mathrm{act}}\) be the nonzero Jacobi support in
the cusp convention of Chapter~\ref{ch:view1-automorphic}, for the
index-one form of \cite{EZ:1}. A finite
set can be obtained by taking
\[
 A_R^\circ=
 \{(n,l,m)\in\Gamma_{\mathrm{act}}:n+|l|+m\le R\}
\]
and its downward saturation in the positive root order. This gives
an exhaustion of the active target degrees. The actual positive
cone and its finite height fibres establish finiteness; they do not
supply geometric representatives.

For geometric charge data, the normal-ordered lift is
\[
 \widehat\Gamma_R=
 \{(c,T):c\in\Gamma_R^{HN},\ T\in\mathcal T_R(c)\},
 \qquad
 \Gamma_R^\Pi=\{\Pi_X(c)+T:(c,T)\in\widehat\Gamma_R\}.
\]
The translate set \(\mathcal T_R(c)\), its bounds and the cocycle
\(B\) are those of Definition~\ref{def:finite-target-test-window}.
The formal lift of every Gram triple in
Lemma~\ref{lem:primitive-lift-every-gram-triple} is a lattice
construction. It does not prove effectivity, stability or existence
of a compact object at that lift.

At a retained positive root \(\beta\), the target has separately
specified even and odd dimensions \(e_\beta,o_\beta\).
The minimal Gritsenko--Nikulin presentation and its ordinary and
super denominators \cite{GN,GNPartI,GNII} determine them by
Proposition~\ref{prop:bkm-all-height-parities}.
Signed exponents alone do not: adding an even and an odd generator
can leave their difference unchanged. The first interior target is
\(29\mid93\), and the next height-four timelike target is
\(201\mid93\). Geometric recognition requires maps to these
spaces, not just the numbers \(-64\) and \(108\).

A finite relation diagram enlarges its degree set to retain the
output of every intermediate bracket. Let \(W_3\) be the thirteen
positive root degrees of height at most three, and put
\[
 R_4=\{3\delta_i+\delta_j:i\ne j\},\qquad
 I_4=\{2a_{ij}:i<j\},
\]
\[
 C_{\le7}=\{a_{ij}+s\delta_k:
       \{i,j,k\}=\{1,2,3\},\ 2\le s\le5\}.
\]
Here \(a_{ij}=\delta_i+\delta_j\), as in the target presentation.
The degree set
\[
 W_{\mathrm{rel}}^+
 =W_3\cup R_4\cup I_4\cup C_{\le7}
       \cup\{2\delta_{123}\},
 \qquad
 D_k=\delta_k+2a_{ij},
\]
has \(35\) elements, including the nine zero root spaces at the
terminal real and complementary Serre degrees.
Adding the \(D_k\) gives
\(W_{\mathrm{sat}}^+=W_{\mathrm{rel}}^+\cup\{D_1,D_2,D_3\}\),
but the effect of a further enlargement must be computed from its
actual set and operations. The first further enlargement described
after Proposition~\ref{prop:target-admissibility-obstruction-wrel}
has \(89\) degrees. Later finite totals cannot be inferred
by retaining a count calculated from a different initial set.

The six-degree families \(A_N,B_N\) of
Proposition~\ref{prop:recognition-affine-staircase} satisfy
\(A_N\Rightarrow B_N\Rightarrow A_{N+2}\) under its specified
degree-enlargement rule. To realize such degrees geometrically,
one must give
\(R_N,c_N,T_N\) satisfying
\(\Pi_X(c_N)+T_N=\gamma_N\) for every relevant member of the
family. Definition~\ref{def:relation-closed-target-finite-r-exhaustion}
keeps this existence assertion separate from the target calculation.

\subsection{Finite moduli and universal complexes}
\label{app:build-finite-moduli}

A bounded geometric class can be described by
\[
 c=(\widehat c,[a_c,b_c],(P_{c,i}),N_c),
\]
where \(\widehat c\) is its charge, \([a_c,b_c]\) its prescribed
cohomological range, \(P_{c,i}\) its Hilbert polynomials and
\(N_c\) a regularity bound. A supplied finite collection \(C_R\)
of such data determines the moduli problem
\[
 \mathfrak M^{ss}_{R,c}
 =\{A\in\operatorname{Perf}(X):
   A\text{ has the specified class and bounds and is semistable}\}.
\]
One must prove that the chosen locus is algebraic and of finite
type. Semistability is not generally a closed condition under every
specialization. If a proper correspondence argument uses a closed
locus inside a bounded Quot space, that closedness is an extra
hypothesis on the locus being used.

The tangent complex at \(A\in\operatorname{Perf}(X)\) is
\(R\operatorname{Hom}(A,A)[1]\). On a CY3 it can have cohomology
in degrees \(-1,0,1,2\). Scalar automorphisms and their Serre-dual
trace direction are present before reduction. A two-term reduced
complex therefore needs an actual reduction map and proof of its
amplitude. Restriction of the ambient shifted symplectic form to an
arbitrary substack is not automatically nondegenerate
\cite{PTVV2013}.

On the scalar gerbe, the universal complex \(\mathcal U\) has
weight one. The complex
\(R\mathcal Hom(\mathcal U,\mathcal U)\) has weight zero and
can descend along the gerbe with its natural maps. Descent of
\(\mathcal U\) itself needs a trivialization or compensating
weight, or else a twisted recipient. The distinction is already
visible on \(B\mathbb G_m\): the standard representation does
not descend to a line on a point, although its endomorphism line
does.

Removing the common scalar subgroup does not make all stabilizers
finite. For a direct sum of two distinct stable objects, a relative
\(\mathbb G_m\) remains; repeated summands retain larger linear
groups. A finite-inertia locus and its closed strata must be
specified if the coefficient or descent theorem requires them.
Quotient by translation in \(E\) is a further geometric operation,
with its own stabilizers and support conditions.

The retained object set, HN factors and chosen extension diagrams
must be compatible. This is not closure under all extensions inside
a bounded nonzero positive charge set: repeated sums leave any
fixed positive bound. Likewise, reversal of a finite HN word and a
formal charge involution does not construct derived duality on the
chosen heart. A geometric duality must land in the specified
recipient and preserve the retained support and orientation data.
The hypotheses and their valid finite consequences are separated
in Appendix~\ref{app:stability-support}.

\subsection{Hybrid support and correspondence maps}
\label{app:build-hybrid-wrapped}

The elliptic support degree is defined from the actual one-cycle:
\[
 (p_E)_*\operatorname{cyc}_1(F)=b_R^{\mathrm{geom}}(F)[E].
\]
For effective one-dimensional support, a positive degree projects
onto all of \(E\). At degree zero, every curve component has a
point as its image in \(E\). These support statements are independent of the
normal-ordered Borcherds coordinate
\(m_R^{\mathrm{Bch}}=\operatorname{pr}_3\overline\Pi_X\).
An equality between those two degree functions is a comparison to
be proved, not a convention for the support.

The local closed-configuration stack for a finite index set \(I\)
has pairs \((A,(e_i)_{i\in I})\) with
\[
 \operatorname{Supp}(A)\subset
       \bigcup_{i\in I}S\times\{e_i\}.
\]
Repeated points are allowed in this closed configuration recipient.
The open disjoint locus and the collision locus have different
restriction maps. A local section on one does not itself define
the operation across the other.

The wrapped sector begins with an \(E\)-equivariant prequotient
\(\mathcal M_{\eta,R}^{\mathrm{wr,rig}}\) and its anchor
\(\lambda_{\eta,R}\). A determinant anchor of Euler weight
\(\chi\) transforms under translation with that weight. It can
be constant along orbits when \(\chi=0\). Jordan--H\"older
spectral data may distinguish, for example, \(\mathcal O_E^2\)
from \(L\oplus L^{-1}\), but translation acts trivially on the
isomorphism classes of degree-zero line bundles. Thus that datum
alone does not separate every orbit.

Changing the origin from \(0\) to \(a\) changes the relevant
Picard class by
\[
 \lambda_a\otimes\lambda_0^{-1}
       \simeq\mathcal O_E\bigl(\chi(0-a)\bigr)
\]
under the stated determinant convention. It does not canonically
choose a stabilizer character. Even a trivial line can carry
different linearizations. An equivariant comparison must retain
that choice.

The four binary correspondences \(LL,LW,WL,WW\) are defined
before quotient, keeping both input anchors and the output anchor.
For two wrapped inputs, the common-translation quotient satisfies
\[
 [M\times N/E]\simeq[M/E]\times_{BE}[N/E]
\]
with diagonal translation. The ordinary product of quotient stacks
would permit two unrelated torsors and changes the operation.
Quotienting an extension diagram therefore requires a
pseudofunctor \(Q_{E,R}\), its operation comparisons
\(\theta^Q_{\mu,R}\), and their base-change coherences.
If \(\mathcal B_R^{E,\mathrm{hyb}}\) denotes the specified
prequotient correspondence construction, the descended one is
\(Q_{E,R}\mathcal B_R^{E,\mathrm{hyb}}\) only after those
comparisons are supplied.

The two-step flag atlas has eight words of length three in
\(\{L,W\}^3\). Their associativity squares are finite tests of
the actual correspondences. Full factorization needs all higher
flags and collision coherences as well. A finite table of binary
maps does not provide those higher maps.

On a sheaf locus, properness of the target extension map can follow
from its description as a closed relative Quot locus. The argument
requires the bounded Hilbert data and the closed source condition.
It applies separately to \(LL\), the two ordered mixed cases,
and \(WW\) when their respective hypotheses hold. It supplies no
properness of the other map in the correspondence and does not
prove Thom--Sebastiani, quotient descent or a protected coefficient
identity. Section~\ref{subsec:hybrid-ran-carrier} gives the actual
local models and the remaining global comparisons.

\subsection{Reduced orientation and equivariant descent}
\label{app:build-orientation}

Let \(K_0\) be a specified perfect deformation complex with a
cosection morphism
\(\sigma:K_0\to\mathcal O[-2]\). Its fibre is perfect:
\[
 K^{\mathrm{fib}}\longrightarrow K_0
       \xrightarrow{\sigma}\mathcal O[-2].
\]
The notation \(K^{\mathrm{red}}\) further specifies which scalar,
orbit and dual obstruction directions are removed. A morphism to
\(H^2(S,\mathcal O_S)[-1]\) would remove a different degree.
The precise shift is therefore part of the cosection construction.

For a nonzero elliptic translation direction \(u_F\), Serre
duality pairs \(\operatorname{Ext}^1(F,F)/\mathbb Cu_F\)
with \(\ker\sigma_F\subset\operatorname{Ext}^2(F,F)\).
This gives a linear reduced pairing. A family of reduced critical
charts, its orientations and coefficient gluing must still be
constructed. The positive reduced stable-pair comparison and the
ordinary critical CoHA construction in the preceding chapters
retain their own geometric domains; neither automatically provides
this complete reduced compact coefficient system.

The determinant convention is
\[
 \det K=\bigotimes_i(\det K^i)^{(-1)^i},\qquad
 \det(K[1])=(\det K)^{-1},\qquad
 \det(K^\vee)=(\det K)^{-1}.
\]
A root is the pair \((o,\rho:o^{\otimes2}\simeq\det K)\).
If a further perfect cosection factor \(C_{\mathrm{cs}}\) is
used, root the actual tensor product
\(\mathscr D=\det K^{\mathrm{red}}\otimes\det C_{\mathrm{cs}}\).
Its obstruction is \(\delta(\mathscr D)\in H^2(B,\mu_2)\).
An underlying cokernel sheaf need not have a determinant until a
perfect model is supplied.

For a complex line,
\[
 w_2(\mathscr D_{\mathbb R})=c_1(\mathscr D)\bmod2.
\]
This uses the integral Chern class. A rational Chern-character
calculation cannot detect all its torsion. If a cocycle
\(a^{\mathrm{red}}\) represents the obstruction, an actual
trivialization is a cochain \(h^{\mathrm{red}}\) with
\(\delta h^{\mathrm{red}}=a^{\mathrm{red}}\).
Vanishing of its cohomology class does not select that cochain.
The isomorphism classes of root pairs form an
\(H^1(B,\mu_2)\)-torsor, rather than merely the two-torsion subgroup
of the Picard group.

For an \(E\)-action, solve the root problem on \([B/E]\).
The group \(H^2(BE;\mathbb F_2)\) has basis \(u_1,u_2\), but
an expression \(a_1u_1+a_2u_2\) contributes only through the
actual quotient or spectral-sequence map. That map need not be
injective. For the translation action of \(E\) on itself,
\([E/E]\) is a point. Thus requiring both coefficients to vanish
in \(H^2(BE)\) is not a universal free-action descent criterion.

\subsubsection*{Finite stabilizers}

At a finite stabilizer \(H\), distinguish the equivariant Kummer
class from the character of the chosen root. For
\(H=E[2]\simeq(\mathbb Z/2)^2\), write
\[
 \beta=b_{20}x_1^2+b_{11}x_1x_2+b_{02}x_2^2,
 \qquad |x_i|=1.
\]
Restriction to the three order-two subgroups gives coefficients
\[
 r_1=b_{20},\qquad r_2=b_{02},\qquad
 r_3=b_{20}+b_{11}+b_{02}.
\]
Hence
\[
 b_{20}=r_1,\quad b_{02}=r_2,\quad b_{11}=r_1+r_2+r_3,
 \qquad \beta=0\Longleftrightarrow r_1=r_2=r_3=0.
\]
Indeed, the two coordinate restrictions send \((x_1,x_2)\) to
\((t,0)\) and \((0,t)\), while the diagonal restriction sends it
to \((t,t)\).

For a two-primary subgroup \((\mathbb Z/2^a)^2\) with
\(a\ge2\),
\[
 H^*(B(\mathbb Z/2^a)^2;\mathbb F_2)
 =\Lambda(x_1,x_2)\otimes\mathbb F_2[y_1,y_2],
 \qquad |y_i|=2.
\]
A degree-two class is
\(A_1y_1+A_{12}x_1x_2+A_2y_2\).
The mixed class restricts to zero on every cyclic subgroup: on a
cyclic group of order divisible by four, both degree-one factors
are multiples of an exterior generator, while on an order-two
subgroup their pullbacks from \(\mathbb Z/2^a\) vanish.
The functional
\[
 \pi_{12}(y_1)=\pi_{12}(y_2)=0,\qquad
 \pi_{12}(x_1x_2)=1
\]
detects the missing coefficient \(A_{12}\).
For \(E[N]\), the exterior-generator case requires \(4\mid N\);
when its two-primary exponent is two, the polynomial
\(E[2]\) calculation applies instead.
Odd-order groups have no positive-degree mod-two group cohomology,
by averaging.

Even when the degree-two class vanishes, a chosen equivariant root
can have a nontrivial character
\(\lambda^H\in H^1(BH;\mathbb F_2)\).
For \(E[2]\), write
\(\lambda^H=\lambda_1x_1+\lambda_2x_2\); its three restriction
values are \(\lambda_1,\lambda_2,\lambda_1+\lambda_2\).
The pair \((1,0)\) gives a nontrivial character despite a zero
Kummer class. Coarse descent requires that every stabilizer act
trivially on the root fibre. Stack descent and coarse descent are
therefore different requirements.

For an inclusion \(K\hookrightarrow H\), chosen cocycle
primitives require comparisons of the form
\[
 \operatorname{res}h^H-h^K=\delta q^{K,H}.
\]
The comparisons must satisfy their own compatibility for chains of
subgroups and for overlap and stage maps. Cohomological vanishing
at each stabilizer separately supplies no such coherent choice.

\subsubsection*{Cross determinants and critical coefficients}

Write \(\mathscr D_A^{\mathrm{full}}=\det R\operatorname{Hom}(A,A)\).
The block decomposition of this full, unshifted complex gives
\[
 \begin{split}
 \mathscr D_{A\oplus B}^{\mathrm{full}}\simeq{}&
 \mathscr D_A^{\mathrm{full}}\otimes\mathscr D_B^{\mathrm{full}}\otimes
 \det R\operatorname{Hom}(A,B)\otimes
 \det R\operatorname{Hom}(B,A).
 \end{split}
\]
For a reduced complex, one must first identify its corresponding
determinant triangle and any additional reduction factors. When
the displayed cross factor is the one in that triangle, the
orientation comparison requires
\[
 o_{A\oplus B}\simeq o_A\otimes o_B\otimes q_{A,B},
 \qquad
 q_{A,B}^{\otimes2}\simeq
 \det R\operatorname{Hom}(A,B)\otimes
 \det R\operatorname{Hom}(B,A).
\]
For a filtered extension, the corresponding determinant triangle
and its coherent higher-flag identifications replace the split
formula. Omitting \(q_{A,B}\) generally changes the square map.

The critical coefficient \(\Phi_o\) is a constructible object
formed using the orientation and its sign torsor
\cite{BBDJS2015,JoyceUpmeier,JoyceTanakaUpmeier}. It is not the
ordinary tensor \(\Phi\otimes o\) of a constructible complex
with a coherent line. A coefficient map on an extension
correspondence must be an actual morphism in its constructible or
mixed-Hodge recipient, with the required shifts and supports.
The line-level determinant identity alone supplies no such morphism.
Thom--Sebastiani comparisons and their full flag associativity must
be checked for these coefficient maps.

\subsubsection*{Weyl cocycles and their compatible primitives}

At a finite action groupoid \(\mathcal G_R\), let the lifted
arrows have sign defect \(c_{o,R}\). With a fixed finite cochain
model, its trivialization is the linear equation
\[
 d_R^1b_R=v(c_{o,R})
\]
over \(\mathbb F_2\). It has a solution exactly when
\[
 \operatorname{rank}d_R^1
 =\operatorname{rank}[d_R^1\mid v(c_{o,R})],
\]
provided the represented defect is actually a cocycle in that
cochain complex. The matrix ranks test solvability; the vector
\(b_R\) is the chosen trivialization.

Transition maps must act on the cochain models, not only their
cohomology. If
\(\rho^i_{RR'}:C^i(\mathcal G_{R'})\to C^i(\mathcal G_R)\),
require
\[
 d_R^1\rho^1_{RR'}=\rho^2_{RR'}d_{R'}^1,
 \quad \rho^2_{RR'}v(c_{o,R'})=v(c_{o,R}),
 \quad \rho^1_{RR'}b_{R'}=b_R.
\]
A tower of solvable equations need not admit the last compatible
choice without an inverse-system argument.

Suppose \(\eta_{R,S}\) trivializes the quotient orientation on a
stratum \(S\), and \(\Theta_{i,R,S}\) is its chosen transport.
The remaining torsor defect is represented by
\[
 \Delta^\omega_{i,R,S}
 =\Theta^1_{i,R,S}\eta_{R,S}-\eta_{R,s_iS}.
\]
It must be closed, have zero class in the actual degree-one
cohomology, and have a compatible primitive when a chain-level
identification is required. The reduced degree-two defects require
the corresponding cochain comparisons as well.
For identified finite stabilizers, the matrices on their cohomology
must intertwine both classes and characters:
\[
 M_i^H v(\beta^H_S)=v(\beta^{H'}_{s_iS}),\qquad
 N_i^H v(\lambda^H_S)=v(\lambda^{H'}_{s_iS}).
\]
The same assertion for any \(BE\)-component uses its actual
induced map \(M_i^E\).

Once coherent lifts have constant residual generator values
\(\sigma_i^o\), their character on a word is
\[
 \epsilon_o(s_{i_1}\cdots s_{i_m})=
       \prod_{a=1}^m\sigma_{i_a}^o.
\]
Section-compatible skew crossings give \(\sigma_i^o=-1\), hence
\(\epsilon_o=\det_W\), as in
Appendix~\ref{app:pfaffian-wall-normal-form}.
The automorphic homogeneous character is trivial on the chamber
symmetry \(S_3\). Extending a geometric \(W\)-action to the full
semidirect product requires actual chamber-symmetry correspondences
and their lifts. The two scalar extensions
\[
 \widetilde\epsilon_\chi(w,\sigma)=\epsilon_o(w)\chi(\sigma),
 \qquad \chi\in\{1,\operatorname{sgn}\},
\]
show that the \(W\)-restriction alone cannot choose that extension.

At another divisor, a coefficient local system \(\mathcal K_R\)
on the complement has monodromy
\[
 T_{R,S,\eta}=\rho_{\mathcal K_R}(\gamma_{R,S,\eta})
\]
around a specified meridian \(\gamma_{R,S,\eta}\). This requires
local constancy and the stated representation of the fundamental
group. A coherent root line has no such representation until a
flat structure is supplied, and an arbitrary constructible critical
complex need not be locally constant on the proposed complement.
Monodromy at an imaginary-root divisor need not be a Weyl reflection.

Parity also does not choose a root-line action. The Berezinian of a
super vector space and a chosen square root with its equivariance
are different data. Similarly, a duality
\(K_\gamma\simeq K_{-\gamma}^{\vee}[-3]\), when actually
supplied, gives
\[
 \det K_\gamma\simeq\det K_{-\gamma}.
\]
The dual and the odd shift each invert the determinant, so their
effects cancel. A claim about opposite root lines needs the chosen
square-map and duality comparison; it cannot be inferred by
inserting an inverse determinant or an unexplained sign line.

Finally, a compatible family of root pairs is first an object of
the limit of the relevant Picard groupoids. Its effectivity as a
line on an infinite geometric recipient needs an actual comparison
\[
 \operatorname{Pic}^{or}(B_\infty)
       \longrightarrow\varprojlim_R\operatorname{Pic}^{or}(B_R).
\]
An equivalence here is stronger than a statement about individual
classes in \(\pi_0\). Taking isomorphism classes can lose the
coherent isomorphisms and their automorphisms.

\subsection{Hall maps, primitives and pairings}
\label{app:build-hall}

Let \(\mathcal H_{R,c}\) be the specified coefficient realization
of the oriented critical object \(\Phi_{o,R,c}\) on the retained
moduli stack. This can be compactly supported critical cohomology
or a declared Borel--Moore model; its grading, duals and shifts are
part of the definition. Do not identify these realizations by
notation alone.

An extension correspondence has maps
\[
 \mathfrak M_a\times\mathfrak M_b
       \xleftarrow{p}\mathfrak E_{a,b}\xrightarrow{q}
       \mathfrak M_{a+b}.
\]
For a product, specify the pullback or Gysin map appropriate to
\(p\), the critical coefficient morphism \(\theta_{a,b}\),
and the proper or compact-support pushforward appropriate to
\(q\). The notation
\[
 m_{a,b}=q_!\theta_{a,b}p^\dagger
\]
records these chosen maps; \(p^\dagger\) includes its degree
shift. Ordinary pullback on sheaves does not automatically define
a map on compactly supported cohomology. The coefficient and
support comparisons must supply that map in the chosen realization.

For a separated finite-type map with a compactification
\(q=\overline q\,j\), the sheaf-theoretic expression is
\[
 q_!K=\overline q_*j_!K.
\]
A common proper refinement of two compactifications identifies
these objects through the standard proper base-change comparison.
This proves independence of the compactification at the functorial
level; it does not manufacture a missing morphism from the resulting
coefficient object to the desired target coefficient.

On the three-input flag correspondence, associativity compares
\[
 m_{a+b,c}(m_{a,b}\otimes1)
       \quad\text{and}\quad
 m_{a,b+c}(1\otimes m_{b,c}).
\]
The geometric proof must identify the two composed diagrams,
their coefficients and shifts. Higher flags give coherence among
these identifications. Likewise a coproduct requires a separate
coefficient map and admissible pushforward in the reverse diagram.
It is not obtained by formally inverting \(\theta_{a,b}\) unless
that map is an equivalence on the stated recipient.

For a bialgebra comparison, the four-block flag decomposes by
\[
 \mathsf S(a,b;u,v)=
 \{(a_1,a_2,b_1,b_2):
   a=a_1+a_2,\ b=b_1+b_2,
   u=a_1+b_1,\ v=a_2+b_2\}.
\]
With \(M\) and \(D\) denoting the product and coproduct matrices,
its required identity is
\[
 \begin{split}
 D^{a+b;u,v}M^{a,b}
 =\sum_{\mathsf S(a,b;u,v)}
 &(M^{a_1,b_1}\otimes M^{a_2,b_2})
 (1\otimes\tau\otimes1)\\
 &\hspace{-1em}\cdot
 (D^{a;a_1,a_2}\otimes D^{b;b_1,b_2}).
 \end{split}
\]
The braiding \(\tau\) must be the one belonging to the actual
tensor. An ordinary Koszul equation cannot be substituted for a
meromorphic braided equation. Unit and counit diagrams are
additional parts of a bialgebra structure.

For a connected positively graded ordinary bialgebra, the antipode
recursion is
\[
 S(1)=1,\qquad S(x)=-x-\sum S(x')x'',
 \qquad\overline\Delta x=\sum x'\otimes x''.
\]
Positive degree makes it finite on each homogeneous element.
The convolution inverse identities are
\(m(S\otimes1)\Delta=m(1\otimes S)\Delta=\eta\epsilon\).
This antipode preserves total charge; it is not a geometric
charge-reversing duality.

Let \(I=\ker\epsilon\). On this augmentation ideal,
\[
 \overline\Delta(x)=\Delta x-1\otimes x-x\otimes1,
 \qquad P=\ker(\overline\Delta:I\to I\otimes I).
\]
For the ordinary Koszul tensor, the supercommutator
\( [x,y]=xy-(-1)^{|x||y|}yx \) of primitives is primitive.
Expanding the bialgebra identity cancels the two cross terms.
Associativity also gives the super Jacobi identity
\[
 (-1)^{|x||z|}[x,[y,z]]+
 (-1)^{|y||x|}[y,[z,x]]+
 (-1)^{|z||y|}[z,[x,y]]=0.
\]
Thus \(P\)
is a Lie superalgebra. After a proved charge descent,
\[
 P^\Pi_\gamma=
   \bigoplus_{\overline\Pi_X(\widehat c)=\gamma}P_{\widehat c},
 \qquad
 [P^\Pi_{\gamma,\bar p},P^\Pi_{\eta,\bar q}]
   \subset P^\Pi_{\gamma+\eta,\bar p+\bar q}.
\]
The additivity of the descended degree is a hypothesis on the
operations; the central cocycle alone does not establish it.

A canonical primitive projector requires further assumptions.
For the connected cocommutative ordinary Hopf realization over
characteristic zero, the Eulerian logarithm gives the projector
proved in Chapter~\ref{ch:view2-bkm}. A chosen linear splitting
of \(P\hookrightarrow I\) is not that projector, and its
naturality under transition maps must be checked independently.

\subsubsection*{Coherent charge descent}

An additive charge map regrades supplied operations. A chain-level
repair of a nonadditive comparison asks for more: a specified
deformation complex, the cochain representing the defect, and its
null-homotopy. If that recipient is a Hochschild cochain complex,
the proposed equation has the form
\[
 d_{\mathrm{Hoch}}\Theta_\Pi=B_{\mathrm{ch}},
 \qquad |\Theta_\Pi|=1,\quad |B_{\mathrm{ch}}|=2.
\]
Here \(B_{\mathrm{ch}}\) must first be constructed as a closed
cochain from the geometric charge defect. The equation is soluble
exactly when its class vanishes in that specified complex.
Triple homotopies, cyclic refinements, Jacobi compatibility and
radical preservation are further equations on the chosen solution.
They may be denoted
\(\mathfrak o_\Theta^{\mathrm{top}},
\mathfrak o_\Theta^{\mathrm{Hoch}},
\mathfrak o_\Theta^{\mathrm{tri}},
\mathfrak o_\Theta^{\mathrm{Jac}},\mathfrak o_F,
\mathfrak o_\Theta^{\mathrm{cyc}},
\mathfrak o_\Theta^{\mathrm{rad}}\): these names record the
carrier, Hochschild, triple, Jacobi, invariant-form, cyclic and
radical requirements, not seven classes in an unspecified group.

The Loday--Quillen--Tsygan comparison concerns stable matrix Lie
homology and cyclic homology
\cite{LodayQuillen1984,Tsygan1983,LodayCyclicHomology}. Applying
it to this problem requires a map from the actual compact
coefficient system to its recipient. Neither that stable theorem
nor the critical Hall construction of \cite{KSCoHA} supplies
\(\Theta_\Pi\), a reduced orientation or all the listed
operation homotopies.

\subsubsection*{Invariant Lie forms and a separate Hopf pairing}

An even opposite-degree invariant Lie form has matrices
\[
 B(e_{\gamma,\bar p,a},e_{-\gamma,\bar p,b})
       =G_{\gamma,\bar p;ab},\qquad
 G_{\gamma,\bar p}=(-1)^{\bar p}G_{-\gamma,\bar p}^{t}.
\]
Other degree or parity pairings vanish. Invariance is
\[
 B([x,y],z)=B(x,[y,z]),
\]
or equivalently
\(B([x,y],z)+(-1)^{|x||y|}B(y,[x,z])=0\).
If this form is defined through a geometric product and a
charge-zero functional, both that functional and the coefficient
pairing must be constructed. A formal integral symbol does not
supply them.

For supplied homogeneous bases and dual output functionals, the
product and pairing entries are then
\[
 M_{ab}^{c}=\langle e^c,q_!\theta p^\dagger(e_a\otimes e_b)\rangle,
 \qquad
 G_{\gamma;ab}=\operatorname{tr}_0
       m(e_{\gamma,a},e_{-\gamma,b}).
\]
The first expression is the coefficient realization of the
correspondence; the second additionally requires its charge-zero
trace and the claimed invariant-form identity. Integration and
duality must be defined on those same coefficient objects.

The radical \(R_B\) is a Lie ideal by invariance. Its generated
ideal in \(U(P)\) is a Hopf ideal because every radical generator
is primitive. This is the positive quotient construction of
Proposition~\ref{prop:bkm-lie-radical-quotient}.
If matrices \(Q_\gamma\) have columns lifting bases of the
radical quotients, then the descended form is
\[
 \overline G_\gamma=Q_\gamma^tG_\gamma Q_{-\gamma}.
\]
The left and right radicals must be killed before this expression
is independent of the lifts. Nondegeneracy requires
\[
 \operatorname{rank}\overline G_\gamma
       =\dim(P_\gamma/R_{B,\gamma})
       =\dim(P_{-\gamma}/R_{B,-\gamma}).
\]
The \(Q\)'s here are sections represented by columns, not arbitrary
quotient matrices denoted by the same symbol.

A full Hopf pairing, if supplied on a different appropriate
recipient, has the additional adjointness equation
\[
 \langle xy,z\rangle
       =\langle x\otimes y,\Delta z\rangle.
\]
For homogeneous basis elements it reads
\[
 \sum_c M_{ab}^{c}G_{cd}
 =\sum_{u,v}D_d^{uv}(-1)^{|e_b||e_u|}G_{au}G_{bv},
\]
together with its unit and counit equations. This is not a further
identity of the invariant Lie form on an ordinary nonabelian
primitive algebra. Indeed, Hopf adjointness and
\(\langle1,x\rangle=\epsilon(x)\) make
\(\langle[x,y],z\rangle=0\) for primitives, whereas the
normalized invariant form on \(\mathfrak{sl}_2\) gives
\(B([h,e],f)=2B(e,f)=2\). Those two requirements cannot be
imposed on the same form in that recipient.

\subsubsection*{Matrix transport}

For a transition \(T_\gamma:H_{R,\gamma}\to H_{R',\gamma}\),
product and coproduct preservation are the equations
\[
 \begin{split}
 T_{\alpha+\beta}M_R^{\alpha,\beta}
   &=M_{R'}^{\alpha,\beta}(T_\alpha\otimes T_\beta),\\
 (T_\mu\otimes T_\nu)D_R^{\rho;\mu,\nu}
   &=D_{R'}^{\rho;\mu,\nu}T_\rho,
 \qquad \rho=\mu+\nu.
 \end{split}
\]
Preserving unit and counit then intertwines the reduced coproduct
on augmentation ideals and hence maps primitives to primitives.
Radical preservation also needs the actual pairing comparison.
For example, if an opposite-degree map \(L\) satisfies
\[
 B_{R'}(T_\gamma x,y')=B_R(x,L_{-\gamma}y'),
\]
then a radical vector maps to a radical vector provided the displayed
pairings test all target vectors. This is a statement about the
specified maps; equal matrix ranks alone do not imply it.

The super bracket matrix is obtained from the two ordered products
and their Koszul signs. Kernel bases, quotient sections and the
maps \(A_{\beta,\bar p}\) to the target must satisfy these
transport equations. The relation complex and its spectral sequence
must use the same representatives and maps. PBW comparisons concern
the actual filtered enveloping algebras and their finite weighted
pieces, not a formal symmetric algebra on a root fragment.

A Hall--Drinfeld double additionally needs a compatible opposite
half, a full pairing in its declared tensor, cross relations and
completion. Serre duality of the geometric complexes and the
weight-five denominator do not supply those constructions
\cite{CYQGVolIII}.

\subsection{Finite skew blocks and the modular section}
\label{app:build-dirac-pfaffian}

For a bundle \(V\) of rank \(2n\) and a skew map
\(D:V\to V^\vee\), the complex
\([V\xrightarrow D V^\vee]\) in degrees \([-1,0]\) has
Pfaffian line \(\det V^\vee\) and section
\(\operatorname{pf}(D)\), with square \(\det D\).
A geometric comparison must include this shift and skew structure;
a fibrewise Serre-dual pairing alone does not supply a family of
skew operators or a section.

For one positive degree \(\gamma\), the block
\[
 J_\gamma=
 \begin{pmatrix}0&1-x^\gamma\\-(1-x^\gamma)&0\end{pmatrix}
\]
has Pfaffian \(1-x^\gamma\). With \(e_\gamma\) even and
\(o_\gamma\) odd copies, the formal super ratio is
\[
 (1-x^\gamma)^{e_\gamma-o_\gamma}.
\]
At finite height, the product of these blocks gives a finite jet.
Local finiteness in positive height makes the jets compatible. It
constructs a formal reference product, not an independently defined
compact geometric skew complex.

For geometric Pfaffians, require a parameter map to the relevant
modular domain and line isomorphisms
\[
 \iota_{\mathrm{aut},R}:\mathscr L_{\mathrm{Pf},R}
       \xrightarrow{\sim}
       (\lambda_{\mathrm{aut}}^5\otimes\varepsilon_5)|_R,
\]
commuting with square maps, stage transitions and cusp frames.
In those frames the desired finite expansion is
\[
 64q^{1/2}r^{1/2}s^{1/2}
       \prod_{\gamma\in A_R}(1-x^\gamma)^{f_\gamma}
       \pmod{\text{higher positive height}},
\]
with the entire boundary factor included. If a protected integration
formula is used to obtain this section, it must act on the actual
coefficient objects and line-valued sections. Defining every input
block to be the desired scalar factor does not construct that
geometric integration map.

The normal wall block is \(D_u=\left(\begin{smallmatrix}0&u\\-u&0\end{smallmatrix}\right)\).
An invariant tangential unit and a section-compatible reflection
lift give the residual sign \(-1\). The global atlas and root
transport are the separate requirements of
Appendix~\ref{app:pfaffian-wall-normal-form}.
The quotient \(L_{\mathrm{II}}/2L_{\mathrm{II}}\) has eight
elements, while the abstract twelve-sign set has \(4096\).
Neither supplies a geometric wall-atlas count without its comparison.

The reference normalization is
\[
 [q^{1/2}r^{1/2}s^{1/2}]\Delta_5=64,\qquad
 [qrs]\Delta_5^2=4096,\qquad
 D_5=64^{-1}\Delta_5,\quad\Xi=D_5^2.
\]
The homogeneous character on the type-II Weyl group is \(\det_W\)
and is trivial on chamber symmetries. The symplectic multiplier is
kept with the weight-five automorphy factor. On the period cover,
\(\operatorname{div}(\Delta_5)=\widetilde H_1\), so its inverse
square has pole order two; the generic coarse quotient has order
one, as in Chapter~\ref{ch:view4-mirror}. These target data do not
construct the geometric section or its mirror interpretation.

Analytic convergence in a cusp region, continuation and boundary
extension are further requirements beyond the formal product.
Under the independent geometric line and section comparisons, the
finite Pfaffian formula yields \(\Delta_5\) at that normalization.
An identity for a protected trace additionally requires its own
state and sewing map; the scalar normalization alone is insufficient.

\subsection{The associative source coalgebra}
\label{app:build-koszul}

Fix the coefficient category, its unit and convolution tensor.
The source \(C_{X,R}\) is a counital coaugmented coalgebra with
differential and a conilpotent or specified positive filtration.
Its cobar is the associative algebra \(\Omega C_{X,R}\).
Conversely an augmented associative source \(A_{X,R}\) has a bar
coalgebra \(BA_{X,R}\). A twisting cochain and the bar--cobar
comparison preserve these types. A coalgebra is not identified
with an endomorphism algebra merely because their underlying
complexes have the same graded ranks.

For an ordinary augmented dg algebra \(A\) over a field \(k\), the underlying bar
coalgebra is the tensor coalgebra on the suspension of its
augmentation ideal. Its length filtration is
\[
 F_p BA=k\oplus\bigoplus_{1\le j\le p}
                (s\overline A)^{\otimes j},
\]
and deconcatenation is
\[
 \Delta[a_1|\cdots|a_p]
 =1\otimes[a_1|\cdots|a_p]+[a_1|\cdots|a_p]\otimes1
  +\sum_{i=1}^{p-1}[a_1|\cdots|a_i]\otimes[a_{i+1}|\cdots|a_p].
\]
Internal differentials and adjacent products give its bar
differential with the suspended Koszul signs. If every factor
has positive height, a bounded total-height piece has bounded
word length. This does not identify word cuts with every
geometric collision operation.

Given coalgebra, source-counit and algebra maps
\[
 b_R:C_{X,R}\longrightarrow BA_{X,R},\qquad
 \varepsilon_R:\Omega BA_{X,R}\longrightarrow A_{X,R},\qquad
 \vartheta_R:A_{X,R}\longrightarrow\mathcal A_{\Delta,E,R},
\]
their composite is
\(\Theta_{\mathrm{Kos},R}=\vartheta_R\varepsilon_R\Omega b_R\).
Its equivalence requires the hypotheses and inverses for these
actual maps. A source-internal counit cannot replace the independent
map \(\vartheta_R\) to the target. The same notation in a Ran
recipient uses its specified convolution and completed bar.

On the positive Ran coefficient recipient of
Section~\ref{subsec:koszul-source}, a proposed comparison has the
form
\[
 C_{X,R}\longrightarrow C_{\Delta,R},\qquad
 \Omega C_{X,R}\longrightarrow\mathcal A_{\Delta,E,R},
\]
with the actual differentials, supports, completions and inverse
maps. The ordinary bar of the reference enveloping algebra is a
separate object. An envelope of a Beilinson--Drinfeld current Lie
algebra and its main-diagonal Ran object also require their specified
functors and units. Francis--Gaitsgory Lie--commutative duality
\cite{FrancisGaitsgoryCKD} does not identify these associative
constructions; the chiral operations of \cite{BD} retain their
specified tensor and support as well.

The finite comparison must intertwine the geometric primitive map,
Weyl transport, root-pair and Pfaffian comparisons and the invariant
Lie form wherever these structures occur. On ordinary enveloping
algebras it must preserve the PBW filtration. Additional full Hopf
pairings, collision maps and factorization operations have their own
comparison equations. Commuting finite maps give a comparison of
pro-objects; passing to an effective geometric or analytic limit
uses separate hypotheses.

\subsection{Recognition from geometric representatives}
\label{app:build-recognition}

Let \(\mathfrak F\) be the free Lie superalgebra on the target
simple and Cartan symbols, and let \(J_{\mathrm{BK}}\) be the
Borcherds relation ideal. Geometric representatives define a map
\[
 \pi:\mathfrak F\longrightarrow P_X^{\Pi,\mathrm{red}}.
\]
Its factorization through the target requires every Cartan,
Chevalley, real Serre, orthogonality and super sign relation.
The desired identification follows when
\[
 \ker\pi=J_{\mathrm{BK}},\qquad
 \operatorname{im}\pi=P_X^{\Pi,\mathrm{red}},
\]
with the stated degree, parity and invariant-form compatibilities.
This is Theorem~\ref{thm:primitive-recognition}.

A finite proof retains bounded word spaces and their actual
intermediate degrees. Its matrices include the product \(M\),
coproduct \(D\), unit and counit, bracket \(B\), invariant form
\(G\), radical inclusion \(K\), quotient section \(Q\), and
comparison maps \(A_{\beta,\bar p}\). The same symbols must not
hide different domains: \(B\) here is a bracket matrix, whereas
the charge normal-ordering \(B(c,c')\) is a cocycle.
The Hall boundary complex, its spectral sequence and the relation
kernel calculations are performed on the chosen geometric
representatives, not on an independently manufactured target vector
space.

For the target positive roots, the height windows
\[
 W_\nu=\{\beta\in\mathcal R_+:
                    \operatorname{ht}\beta\le\nu\},\qquad
 \bigcup_\nu W_\nu=\mathcal R_+,
\]
give the Gram degree
\(\gamma_\beta=(b_1,b_1+b_2-b_3,b_2)\) for
\(\beta=b_1\delta_1+b_2\delta_2+b_3\delta_3\).
Thus \(\Gamma(W_\nu)=\{\gamma_\beta:\beta\in W_\nu\}\)
is a target degree window. Its exhaustion is proved by
\(\beta\in W_{\operatorname{ht}\beta}\); source support
exhaustion requires the separate geometric representatives.

The finite conditions \textup{(R0)}--\textup{(R8)} in
Definition~\ref{def:cofinal-primitive-recognition-datum} retain
source supports, representatives, relation identities, parity
bijections, the invariant form and radical, kernel equality,
generation, PBW and transition compatibility. Every relation word
must be evaluated in a stage containing its intermediate outputs.
The all-level assertion follows only when these bounded word spaces
exhaust the presentation with compatible comparisons. A finite
prefix or a list of signed multiplicities does not prove it.

At \(\delta_{123}\), recognition compares \(29\) even and
\(93\) odd vectors together with their brackets. The rank-28
differential of Chapter~\ref{ch:view2-bkm} extends through height
five in the stated presentation and fails at height six. A derived
replacement changes its presentation or its recipient and must
provide the corresponding map. It cannot silently keep the original
relations while retaining the obstructed differential.

For a topological conclusion, specify the completion of both
source and target. Positive-height completion uses genuine
positive Lie quotients. A product over both positive and negative
root degrees does not automatically have continuous brackets:
a fixed output degree can receive infinitely many cancelling input
pairs. A different completion requires its own finite-support or
convergence argument.

An ordinary primitive isomorphism induces the filtered enveloping
Hopf isomorphism under its stated hypotheses. It does not by itself
identify a chiral field algebra, internal centre, boundary category
or physical state. Those are further operation-preserving maps
\cite{ChiralBarCobarVolI,ChiralBarCobarVolII,CYQGVolIII}.

\subsection{Dependencies, transitions and effectivity}
\label{app:build-ordered-dependencies}

The geometric dependencies form a partial order. The Hall operations
need their correspondence and coefficient maps; recognition needs
the actual source primitives and their presentation map. A geometric
Pfaffian comparison needs a skew family, root data and a section.
These are dependencies of this proposed compact construction, not
a universal forced order among all the objects. A determinant root
can be studied before any Hall product, and the bar of an augmented
algebra is defined without an orientation line.

The finite morphisms \textup{(M1)}--\textup{(M8)} of
Definition~\ref{def:finite-stage-dirac-igusa-morphism} express the
required compatibility. They must preserve charge, coefficients,
operations, root pairs, Pfaffian sections, coalgebras and recognition
maps. Root-line transitions are isomorphisms after the appropriate
pullback. Kernel and pairing transitions use their actual source
and target maps. Cohomological ranks alone do not define them.

The collection \(\{\mathfrak D^{\mathrm{DI}}_{X,R}\}_R\) is a
pro-object once the maps and coherence are supplied. Its Hom formula
is the limit-colimit formula of
Proposition~\ref{prop:dirac-igusa-pro-object-universal-property}.
Compatible levelwise inverses are sufficient for a pro-isomorphism,
but not necessary in a general pro-category. The zero-bonding
nonzero-stage tower in that subsection gives a counterexample to
the converse.

For each countable coefficient tower, the exact obstruction to
commuting cohomology with derived inverse limit is
\(\varprojlim{}^1 H^{q-1}\) \cite[\S3.5]{WeibelHA}.
Finite-dimensional field-valued
towers are Mittag--Leffler; finite-rank integral towers need not be.
Primitive-space, pairing-kernel and PBW towers require this check
on their own coefficients. A nonuniform collection of finite
cutoffs gives no all-level conclusion without a cofinality argument.

The resulting conditional comparison is the one in
Theorem~\ref{thm:pfaffian-dirac-conditional}. It retains the
geometric degeneration and coefficient, quotient orientation,
coherent Weyl transport and wall-atlas conditions, together with
independent Pfaffian, primitive and associative comparison data.
Its remaining compact construction is to supply these objects and
maps on the same geometric system. Their formal reference models,
finite matrix identities and automorphic scalar are necessary
comparison data, but do not replace that construction.
